% ============================================================================
\chapter{Composer, standards PSR, qualit\'e et outils de development}
\label{ch:php-composer}
% ============================================================================

\begin{objectif}
\begin{itemize}
  \item Utiliser Composer comme un adulte~: manifeste, lock, scripts, paquets, publication.
  \item Conna\^itre les PSR qui changent r\'eellement le code (1, 3, 4, 6, 7, 11, 12, 15, 17).
  \item Automatiser le style et la correction~: PHP-CS-Fixer/PHP~Codesniffer, PHPStan, Rector.
  \item Installer l'environnement serveur~: nginx + php-fpm, OPcache, Xdebug, journaux.
  \item Tester~: PHPUnit, Mockery, DBunit-like, et les patterns pour rendre le code testable.
\end{itemize}
\end{objectif}

\section{Composer}

\begin{terminal}
# Installation officielle (Linux/macOS)
curl -sS https://getcomposer.org/installer | php -- --install-dir=/usr/local/bin --filename=composer
composer --version
# Windows : https://getcomposer.org/Composer-Setup.exe
# Alt : php -r "copy('https://getcomposer.org/installer', 'composer-setup.php');" ... (verifie la signature)

cd mon-projet
composer init
composer require guzzlehttp/guzzle monolog/monolog          # prod
composer require --dev phpunit/phpunit phpstan/phpstan rector/rector friendsofphp/php-cs-fixer
composer update                              # respecte les contraintes, puis reecrit composer.lock
composer update vendor/paquet --with-all-dependencies
composer outdated --direct                   # uniquement les dependances directes
composer show --tree
composer why symfony/polyfill-ctype          # qui tire ca, et pourquoi
composer validate --strict
composer dump-autoload -o --classmap-authoritative    # en prod : table classe->fichier, pas de stat
composer audit                               # CVE connues dans l'arbre
composer check-platform-reqs
composer install --no-dev --optimize-autoloader        # la commande de deploiement
\end{terminal}

\begin{codejson}
{
  "name": "nour/blog",
  "description": "Un blog PHP moderne, sans framework.",
  "type": "project",
  "license": "MIT",
  "require": {
    "php": ">=8.3",
    "ext-pdo": "*",
    "ext-mbstring": "*",
    "ext-json": "*",
    "monolog/monolog": "^3.5",
    "symfony/dotenv": "^7.0",
    "vlucas/phpdotenv": "^5.6"
  },
  "require-dev": {
    "phpunit/phpunit": "^11.0",
    "phpstan/phpstan": "^1.11",
    "rector/rector": "^1.2",
    "friendsofphp/php-cs-fixer": "^3.54",
    "squizlabs/php_codesniffer": "^3.9",
    "roave/security-advisories": "dev-latest"
  },
  "autoload": { "psr-4": { "App\\": "src/" } },
  "autoload-dev": { "psr-4": { "Tests\\": "tests/" } },
  "scripts": {
    "test": "phpunit --colors=always",
    "stan": "phpstan analyse --memory-limit=512M",
    "fix": ["php-cs-fixer fix", "rector process"],
    "check": ["@fix", "@stan", "@test"],
    "serve": "php -S 127.0.0.1:8000 -t public public/router.php",
    "post-install-cmd": ["@php bin/console migrate"]
  },
  "scripts-descriptions": { "check": "Formate, analyse, teste." },
  "config": {
    "sort-packages": true,
    "optimize-autoloader": true,
    "allow-plugins": false,
    "platform": { "php": "8.3.0" }
  },
  "minimum-stability": "stable",
  "prefer-stable": true
}
\end{codejson}

\etape{Semver et contraintes---les m\^emes notions qu'en npm, deux pi\`eges en plus}
\begin{itemize}
  \item \verb|^7.0| accepte \verb|7.x| (comme npm). \verb|~7.0| accepte \verb|7.0.x| -- et
        \emph{aussi} \verb|7.1| en philosophie Composer, ce qui surprend~: sois explicite.
  \item \verb|dev-main| et les \verb|@dev| ne se m\'elangent pas avec des versions stables~; \'evite
        les branches stables en production (impossible de reconstruire \`a l'identique sans le lock).
  \item \texttt{composer.lock} se versionne~: pour une \emph{biblioth\`eque} il est ignor\'e, pour une
        \emph{application} il est sacr\'e. C'est la seule r\`egle qui distingue les deux cas.
  \item \verb|platform.php| g\`ele la version \og ~cible~ \fg{}~: utile quand ton poste n'a pas la m\^eme
        version que le serveur (et qu'un paquet exige \og ~php \textgreater=~8.4~ \fg{}).
\end{itemize}

\etape{Publier une biblioth\`eque}
\begin{terminal}
# sur GitHub : composer.json correct, tag sement, README, licence
git tag v1.0.0 && git push --tags
# rendre le paquet decouvrable : https://packagist.org/submit (souvent auto via GitHub Actions)
# chez les autres : composer require nour/blog-utils "^1.0"
# paquet prive : repository "vcs" (https://repo.privie.git) + config gitlab-token,
#   ou Satispress/Private Packagist, ou un depot artefacts.
\end{terminal}

\section{Les standards PSR, en pratique}

\begin{tabz}{llX}
  \headrow \thd{PSR} \& \thd{Nom} \& \thd{Ce que \c{c}a change chez toi} \\
  \midrule
  1/12 \& Coding Style \& \verb|declare(strict\_types=1);|, un fichier = une classe, \verb|namespace App\;|, \verb|use| tri\'es, accolade ouvrante \`a la ligne suivante pour les m\'ethodes, indent 4 espaces. \\
  3 \& Log \& \verb|Psr\\Log\\LoggerInterface| : tu remplaces Monolog par n'importe quoi sans toucher le code. \\
  4 \& Autoload \& un dossier \verb|src/| mapp\'e sur un pr\'efixe d'espace de noms. \\
  6 \& Contexte d'interface \& \verb|getContainer()| etc. : rare en app, fr\'equent en biblioth\`eque. \\
  7/17 \& Messages HTTP \& objets \verb|Request|/\verb|Response| immuables, \verb|ServerRequestInterface| ; PSR-17 les fabrique. \\
  11 \& Conteneur \& \verb|get()|/\verb|has()|~: tous les conteneurs se valent. \\
  14 \& Exceptions \& \verb|ThrowableInterface|~: rare, mais les biblioth\`eques s\'erieuses le font. \\
  15 \& Middlewares \& \verb|MiddlewareInterface|~: le mod\`ele de Slim, Laminas, et des frameworks modernes. \\
  16 \& Cache \& \verb|CacheItemPoolInterface|~: code ton cache une fois, change d'impl\'ementation \`a volont\'e. \\
  18 \& HTTP Client \& \verb|HttpClientInterface|~: d\'ecouple de Guzzle. \\
  20/23 \& Environnement \& erreurs \& lecture d'\verb|.env|~; \verb|bootstrap.php| commun. \\
\end{tabz}

Les liens de lecture~: \link{PHP-FIG}{https://www.php-fig.org/psr/},
\link{PHP~PSR~visual~guide}{https://php-fig.github.io/},
\link{PHP Standard Recommendations sur packagist}{https://packagist.org/providers/psr/log-implementation}.

\section{Style et analyse automatique}

\begin{terminal}
# PHP-CS-Fixer : aligne le style sur un preset (Symfony, PSR-12)
vendor/bin/php-cs-fixer fix --diff --using-cache=no
# PHP_CodeSniffer : verifier sans corriger (id\'eal en CI)
vendor/bin/phpcs --standard=PSR12 --colors src tests
vendor/bin/phpcbf --standard=PSR12 src            # autocorrection
# PHPStan : la vraie valeur ajout\'ee --- trouve les bugs, pas les espaces
vendor/bin/phpstan analyse --level=8 src
# Psalm (analyse plus profonde, plus bavarde)
vendor/bin/psalm --level=2
# Rector : refonte automatique (monter de version PHP, moderniser une API)
vendor/bin/rector process --dry-run
\end{terminal}

\begin{codeconf}
# .php-cs-fixer.dist.php
<?php
$finder = (new PhpCsFixer\Finder())->in(['src', 'tests'])->append(['config']);
return (new PhpCsFixer\Config())
    ->setRiskyAllowed(true)
    ->setRules([
        '@Symfony' => true,
        '@PHP83Migration' => true,
        'declare_strict_types' => true,
        'global_namespace_import' => ['import_classes' => true],
        'ordered_imports' => ['sort_algorithm' => 'alpha'],
        'native_function_invocation' => ['include' => ['@compiler_optimized']],
        'phpdoc_to_comment' => false,
        'yoda_style' => false,
        'no_unused_imports' => true,
        'single_line_throw' => false,
    ])
    ->setFinder($finder);
\end{codeconf}

\begin{codeconf}
# phpstan.neon
parameters:
    level: 8
    paths: [src, tests]
    checkMissingIterableValueType: true
    treatPhpDocTypesAsCertain: false
    reportUnmatchedIgnoredErrors: false
    ignoreErrors:
        - identifier: missingType.iterableValue
    bootstrapFiles: [tests/bootstrap.php]
\end{codeconf}

\begin{note}
Ordre d'adoption qui rapporte~: (1) PHP-CS-Fixer (5 minutes, aucun d\'ebat)~; (2) PHPStan niveau
5, puis 8 progressivement~; (3) \verb|declare(strict\_types=1)| partout~; (4) Rector pour les
mont\'ees de version~; (5) Psalm si le projet est critique. Le niveau 8 de PHPStan est le vrai
d\'etecteur de bugs de nullit\'e, de types fant\^omes et d'oublis de v\'erification---l'\'equivalent du
\verb|strict| de TypeScript, en plus s\'ev\`ere.
\end{note}

\section{Le serveur~: nginx + php-fpm (ou Apache)}

\begin{codeconf}
# /etc/nginx/sites-available/nour
server {
    listen 80;
    server_name nour.dev www.nour.dev;
    root /var/www/nour/public;                 # LA RACINE WEB = public/ , jamais le projet
    index index.php;

    location / {
        try_files $uri $uri/ /index.php$is_args$args;   # le "front controller"
    }
    location ~ \.php$ {
        include fastcgi_params;
        fastcgi_pass unix:/run/php/php8.3-fpm.sock;
        fastcgi_param SCRIPT_FILENAME $realpath_root$fastcgi_script_name;
        fastcgi_param DOCUMENT_ROOT $realpath_root;
        fastcgi_read_timeout 60s;
    }
    location = /favicon.ico { log_not_found off; access_log off; }
    location ~* \.(css|js|woff2|svg|png|jpg|webp)$ {
        expires 1y; add_header Cache-Control "public, immutable"; access_log off;
    }
    location ~ /\.(?!well-known) { deny all; }
}
\end{codeconf}

\begin{codeconf}
; /etc/php/8.3/fpm/pool.d/nour.conf  -- un pool par application
[nour]
user = nour
group = nour
listen = /run/php/php8.3-fpm-nour.sock
listen.owner = www-data
listen.group = www-data
pm = dynamic
pm.max_children = 25                 ; memoire_par_processus ~60 Mo => (RAM*0.8)/60
pm.start_servers = 6
pm.min_spare_servers = 4
pm.max_spare_servers = 10
pm.max_requests = 500                 ; recycle pour limiter les fuites
clear_env = no                        ; sans cela, tes variables d'environnement ne passent PAS
php_admin_value[upload_max_filesize] = 10M
php_admin_value[memory_limit] = 256M
php_admin_value[error_log] = /var/log/php/nour-error.log
catch_workers_output = yes
decorate_workers_output = no
\end{codeconf}

\begin{terminal}
sudo systemctl reload php8.3-fpm nginx
php -i | grep -E "opcache.enable|memory"        # OPcache doit etre actif
# OPcache en production (les memes opcodes recompiles a chaque requete = gaspillage) :
#   opcache.enable=1, opcache.memory_consumption=192, opcache.max_accelerated_files=20000,
#   opcache.validate_timestamps=0  (+ `php-fpm -tt && kill -USR2 <master>` apres un deploiement)
\end{terminal}

\section{D\'ebogueur Xdebug---et comment s'en servir sans pleurer}

\begin{terminal}
# Linux (paquet) : sudo apt install php-xdebug ; ou via le gestionnaire de la distro
# macOS : pecl install xdebug  |  Windows : download.coderix.com/xdebug (choisir le .dll
#         correspondant a php -i : version, NTS/TS, xdebug)
php --ini ; php -i | grep xdebug.mode
\end{terminal}

\begin{codeconf}
; /etc/php/8.3/mods-available/xdebug.ini  (ou dans php.ini)
zend_extension=xdebug
xdebug.mode=debug
xdebug.start_with_request=trigger        ; ne s'active que si XDEBUG_TRIGGER present : + de controle
xdebug.client_host=127.0.0.1
xdebug.client_port=9003
xdebug.idekey=VSCODE
xdebug.log=/tmp/xdebug.log
xdebug.max_nesting_level=512
xdebug.profiler_output_dir=/tmp
xdebug.mode=profile,debug                 ; les deux a la fois en dev local seulement
\end{codeconf}

\begin{codejson}
{
  "name": "Ecouter Xdebug",
  "type": "php",
  "request": "launch",
  "port": 9003,
  "pathMappings": { "/var/www/nour": "${workspaceFolder}" },
  "log": true
}
\end{codejson}

D\'eclencher une session~: \link{xdebug-trigger}{https://xdebug.org/wizard}~; en pratique~: extension
\og ~Xdebug Helper~ \fg{}, ou \cmd{curl -H "XDEBUG-TRIGGER: 1" http://nour.dev/article?id=7}.
\verb|xdebug\_break()| dans le code, \verb|xdebug\_print\_function\_stack()| pour voir qui m'a appel\'e,
\verb|xdebug\_start\_code\_coverage()| pour mesurer ce qui manque.

\subsection{Profilage~: o\`u part le temps}

\begin{terminal}
# 1. Xdebug profile (cachegrind) puis QCachegrind / KCacheGrind / WebCacheGrind
curl -H "XDEBUG-PROFILE: 1" http://nour.dev/rapport -o /dev/null
kcachegrind /tmp/cachegrind.out.*

# 2. blackfire.io (versions locales gratuites) ou xhprof
composer require --dev symfony/xhprof

# 3. Le minitraceur maison, 10 lignes, suffisant 9 fois sur 10
#    $t=microtime(true); ... ; error_log(sprintf('[perf] %s: %.1f ms', $etape, 1e3*(microtime(true)-$t)));
\end{terminal}

\section{Tests PHP~: PHPUnit (et Pest)}

\begin{codephp}
<?php
declare(strict_types=1);
namespace Tests\Unit\Domain;

use App\Domain\Article\Article;
use App\Domain\Article\Statut;
use DateTimeImmutable;
use PHPUnit\Framework\Attributes\CoversClass;
use PHPUnit\Framework\Attributes\DataProvider;
use PHPUnit\Framework\TestCase;

#[CoversClass(Article::class)]
final class ArticleTest extends TestCase
{
    private function article(string $titre = 'Un titre correct', Statut $s = Statut::Brouillon): Article
    {
        return new Article(null, $titre, 'corps', $s, new Auteur('N', 'n@e.com'), new DateTimeImmutable());
    }

    public function testLaPublicationFixeLaDate(): void
    {
        $a = $this->article()->publier();
        self::assertSame(Statut::Publie, $a->statut());
        self::assertNotNull($a->publieLe());
    }

    #[DataProvider('titresInvalides')]
    public function testRefuseLesTitresInvalides(string $titre, string $attendu): void
    {
        $this->expectException(\DomainException::class);
        $this->expectExceptionMessageMatches($attendu);
        $this->article($titre);
    }
    public static function titresInvalides(): iterable
    {
        yield 'vide' => ['', '/minimum/'];
        yield 'trop court' => ['ab', '/minimum/'];
        yield 'trop long' => [str_repeat('a', 200), '/maximum/'];
    }

    public function testRenommerNeMutePasLOriginal(): void
    {
        $original = $this->article('A');
        $modifie = $original->renommer('B');
        self::assertNotSame($original, $modifie);
        self::assertSame('A', $original->titre());
    }
}
\end{codephp}

\begin{codephp}
<?php
// tests/bootstrap.php --- le minimum pour tester sans toucher a la vraie base
require __DIR__ . '/../vendor/autoload.php';
putenv('APP_ENV=test');
// SQLite en memoire pour les tests d'integration : rapide, isole, reel
$pdo = new PDO('sqlite::memory:');
$pdo->exec(file_get_contents(__DIR__ . '/../database/schema.sqlite.sql'));
\end{codephp}

\begin{codeconf}
; phpunit.xml.dist
<?xml version="1.0" encoding="UTF-8"?>
<phpunit xmlns:xsi="http://www.w3.org/2001/XMLSchema-instance"
         xsi:noNamespaceSchemaLocation="vendor/phpunit/phpunit/phpunit.xsd"
         bootstrap="tests/bootstrap.php" colors="true" cacheDirectory=".phpunit.cache"
         failOnWarning="true" failOnNotice="true" beStrictAboutOutputDuringTests="true">
  <testsuites>
    <testsuite name="unit"><directory>tests/Unit</directory></testsuite>
    <testsuite name="integration"><directory>tests/Integration</directory></testsuite>
  </testsuites>
  <source>
    <include><directory>src</directory></include>
  </source>
  <coverage>
    <report><text outputFile="php://stdout"/><html outputDirectory="coverage"/></report>
  </coverage>
  <php>
    <ini name="error_reporting" value="-1"/>
    <server name="APP_ENV" value="test" force="true"/>
  </php>
</phpunit>
\end{codeconf}

\begin{terminal}
vendor/bin/phpunit --testdox --coverage-text
vendor/bin/phpunit --filter=Article --coverage-clover coverage.xml
vendor/bin/phpunit --order-by=defects --reverse-list    # les echecs d'abord
# Pest (syntaxe plus courte, tres appreciee) : composer require --dev pestphp/pest && vendor/bin/pest
\end{terminal}

\subsection{Paquets PHP qui valent leur poids}

\begin{tabz}{lX}
  \headrow \thd{Paquet} \& \thd{Usage} \\
  \midrule
  \link{monolog/monolog}{https://github.com/Seldaek/monolog} \& journalisation PSR-3, handlers fichier/Sentry/Slack.\\
  \link{symfony/dotenv}{https://symfony.com/components/Dotenv} \& lecture de \fichier{.env} par environnements.\\
  \link{nesbot/carbon}{https://carbon.nesbot.com/} \& dates agr\'eables (et \verb|DateTimeImmutable| natif, souvent suffisant).\\
  \link{ramsey/uuid}{https://github.com/ramsey/uuid} \& UUIDv7/v4.\\
  \link{vlucas/phpdotenv}{https://github.com/vlucas/phpdotenv} \& l'alternative simple.\\
  \link{guzzlehttp/guzzle}{https://docs.guzzlephp.org/} ou \link{symfony/http-client}{https://symfony.com/components/HttpClient} \& client HTTP.\\
  \link{phpmailer/phpmailer}{https://github.com/PHPMailer/PHPMailer} ou \link{symfony/mailer}{https://symfony.com/doc/current/mailer.html} \& emails.\\
  \link{respect/validation}{https://respectvalidation.readthedocs.io/} \& validation d\'eclarative.\\
  \link{league/flysystem}{https://flysystem.thephpleague.com/} \& stockage abstrait (local, S3, FTP).\\
  \link{symfony/console}{https://symfony.com/components/Console} \& commandes CLI, sans framework.\\
  \link{robmorgan/phinx}{https://book.cakephp.org/phinx/} \& migrations.\\
  \link{mockery/mockery}{https://docs.mockery.io/} \& mocks (surtout pour les interfaces).\\
  \link{friendsofphp/php-cs-fixer}{https://github.com/PHP-CS-Fixer/PHP-CS-Fixer} \& style.\\
  \link{rector/rector}{https://getrector.com/} \& refactorisations en masse.\\
  \link{filp/whoops}{https://github.com/filp/whoops} \& page d'erreur de d\'eveloppement magnifique.\\
\end{tabz}

\begin{exercice}
\begin{enumerate}
  \item Cr\'ee un projet \og ~sans framework~ \fg{} propre~: \verb|composer init|, PSR-4, \verb|public/index.php|,
        \verb|src/|, \verb|tests/|, \verb|data/|, scripts \verb|serve|, \verb|check|, \verb|test|.
  \item Ajoute PHP-CS-Fixer et PHPStan niveau~8~: fais passer le projet au vert en partant de
        \og ~plein d'erreurs~ \fg{}. Note les bugs r\'eels trouv\'es par PHPStan (il y en a).
  \item \'Ecris 20 tests PHPUnit~: d'abord sur le domaine pur (aucune d\'ependance), puis sur le d\'ep\^ot
        avec SQLite en m\'emoire, puis un test d'int\'egration HTTP (une vraie requ\^ete sur le routeur).
  \item Configure Xdebug~: arr\^ete-toi dans une action, inspecte la pile, utilise
        \verb|xdebug\_break()| dans une fermeture. Puis profile une page lente et nomme la fonction
        qui co\^ute~; corrige, reprofile.
  \item Rends ton d\'ep\^ots compatible PSR-3 et PSR-16~: ajoute un cache m\'emoire avec TTL~; passe
        du \verb|apcu| au \verb|files| sans changer une ligne de m\'etier.
  \item Publie une petite biblioth\`eque sur Packagist~: tests, CI, style, analyse, tag, README~;
        installe-la dans un autre projet.
\end{enumerate}
\end{exercice}

\begin{ressources}
\begin{itemize}
  \item Composer documentation --- \url{https://getcomposer.org/doc/}
  \item PHP-FIG --- \url{https://www.php-fig.org/psr/}
  \item PHPStan --- \url{https://phpstan.org/} --- Psalm --- \url{https://psalm.dev/} --- Rector --- \url{https://getrector.com/}
  \item PHPUnit --- \url{https://docs.phpunit.de/} --- Pest --- \url{https://pestphp.com/}
  \item Xdebug --- \url{https://xdebug.org/docs/step\\\_debug} --- PHP-FPM tuning --- \url{https://php.watch/configuration/php-fpm}
  \item \link{PHP~The~Right~Way}{https://phptherightway.com/} --- relis la partie \og ~Dependencies~ \fg{} et \og ~Tooling~ \fg{}.
\end{itemize}
\end{ressources}

\index{Composer}\index{PSR}\index{PHPStan}\index{Xdebug}\index{PHPUnit}
